\begin{proof}
    %
    We can re-arrange the covariance in the definition of $\infl_n$ as the product
    of terms centered at $\thetahat$:
    %
    \begin{align*}
    %
    \infl_n ={}& N \expect{\post}{
        (\g(\theta) - \g(\thetahat))
        (\ell(\x_n \vert \theta) - \ell(\x_n \vert \thetahat))
    } -
    \\&
    2 N \left(
        \g(\thetahat) - \expect{\post}{\g(\theta)}
    \right)
    \left(
        \ell(\x_n \vert \thetahat) - \expect{\post}{\ell(\x_n \vert \theta)}
    \right).
    %
    \end{align*}
    %
    Apply \thmref{bayes_clt_main} three times, with $\phi_\g(\theta) := \g(\theta)$,
    $\phi_\ell(\theta, \x_n) := \ell(\x_n \vert \theta)$, and $\phi_{\ell\g}(\theta,
    \x_n) := (\g(\theta) - \g(\thetahat)) (\ell(\x_n \vert \theta) - \ell(\x_n \vert
    \thetahat))$, all of which are third-order BCLT-okay by assumption. Let the
    residuals be denoted as $\resid{\g}$ (with no $\x_n$ dependence),
    $\resid{\ell}(\x_n)$, and $\resid{\ell\g}(\x_n)$, respectively.   Note that the
    first derivative of $\theta \mapsto \phi_{\ell\g}(\theta, \x_n)$ is zero at
    $\thetahat$.

    Write the leading term in the expansion of $\phi_\ell(\theta,
    \x_n)$ as
    %
    \begin{align*}
    %
    L(\x_n) :=
    \frac{1}{2} \ellgrad{2}(\thetahat, \x_n) \infohat^{-1} +
    \frac{1}{6} \ellgrad{1}(\thetahat, \x_n) \likhatk{3}(\thetahat) \normfourth.
    %
    \end{align*}
    %
    Under \assuref{bayes_clt}, $\thetahat \plim \thetatrue$  and $\infohat^{-1}$ has
    bounded operator norm with probability approaching one (\lemref{thetahat_consistent}).
    These facts imply that $\meann \norm{L(\x_n)}_2^2 = \ordp{1}$ by Cauchy-Schwarz
    and a ULLN (\assuref{bayes_clt} \itemref{loglik_ulln} and \lemref{ulln}).

    Then \thmref{bayes_clt_main}, together with the fact that $\thetahat = \ordp{1}$
    and $\g(\theta)$ is continuous, gives
    %
    \begin{align*}
    %
    \infl_n ={}&
    \ggrad{1}(\thetahat) \infohat \ellgrad{1}(\x_n \vert \thetahat) +
    \ordlog{N^{-1}} \resid{\ell\g}(\x_n) +
    \\& N
        \left( \ordp{N^{-1}} L(\x_n) + \ordlogp{N^{-2}} \resid{\ell}(\x_n) \right)
        \left( \ordp{N^{-1}} + \ordlogp{N^{-2}} \resid{\g} \right)
    \\=&
    \ggrad{1}(\thetahat) \infohat \ellgrad{1}(\x_n \vert \thetahat) +
    \\&
    % \ordlog{N^{-1}} (\resid{\ell\g}(\x_n) + L(\x_n)) +
    % \ordlogp{N^{-2}} \left(\resid{\g} L(\x_n) +  \resid{\ell}(\x_n)\right) +
    % \ordlogp{N^{-3}} \resid{\ell}(\x_n) \resid{\g}.
    \ordlog{N^{-1}} \left(
    \resid{\ell\g}(\x_n) + L(\x_n) +
    \resid{\g} L(\x_n) +  \resid{\ell}(\x_n) +
    \resid{\ell}(\x_n) \resid{\g}
    \right).
    %
    \end{align*}
    %
    In the final line, we have combined lower orders into the $\ordp{N^{-1}}$ term
    for conciseness of expression.

    We now show that the sums of the squares of the residuals in the preceding
    display go to zero. By \thmref{bayes_clt_main}, each of $\meann
    \resid{\ell}(\x_n)^2$, and $\meann \resid{\ell\g}(\x_n)^2$ are $O_p(1)$.
    Similarly, $\meann \norm{\infohat^{-1} \ellgrad{1}(\x_n \vert \thetahat)}_2^2 =
    \ordp{1}$ by a ULLN and the boundedness of the operator norm of $\infohat^{-1}$.
    By Cauchy-Schwarz, we then have $\meann \ggrad{1}(\thetahat)^\trans
    \infohat^{-1} \ellgrad{1}(\x_n \vert \thetahat) \resid{\infl}(\x_n) = O_p(1)$.

    By repeated application of Cauchy-Schwarz, we
    can thus sum $\infl_n$ and $\infl_n \infl_n^\trans$ to get
    %
    \begin{align*}
    %
    \sumn \infl_n ={}&
     \ggrad{1}(\thetahat) \infohat \sumn \ellgrad{1}(\x_n \vert \thetahat) +
     \ordlogp{N^{-1}} = \ordlogp{N^{-1}} \\
    \meann \infl_n \infl_n^\trans ={}&  \ggrad{1}(\thetahat) \infohat^{-1}
     \scorecovhat
     \infohat^{-1} \ggrad{1}(\thetahat)^\trans + \ordlogp{N^{-1}}
    = \gcovmaphat + \ordlogp{N^{-1}}.
    %
    \end{align*}
    %
    Consistency of $\gcovij$ then follows from \corref{bayes_expectation_dist}.
    %
\end{proof}